% (c)~2014 Dimitrios Vrettos - d.vrettos@gmail.com
% (c)~2014 Claudio Carboncini - claudio.carboncini@gmail.com
% (c)~2015 Daniele Zambelli daniele.zambelli@gmail.com
% (c)~2015 Maria Antonietta Pollini

\input{\folder diseq2gr_grafici.tex}

\chapter{Disequazioni di secondo grado}

\inicapitolo{%
\begin{itemize}[left=0mm, nosep]
\item studio del segno di un polinomio di secondo grado;
\item disequazioni di secondo grado;
\item disequazioni di grado superiore;
\item disequazioni fratte;
\item sistemi di disequazioni.
\end{itemize}
}

\section{Risoluzione delle disequazioni di secondo grado}
\label{sec:diseq_secondo_grado}
\ind{risolvere una disequazione di 2° grado}

Una disequazione di secondo grado si presenta in una delle seguenti forme:
\[ax^2+bx+c < 0;\quad ax^2+bx+c \leqslant 0;
  \quad ax^2+bx+c\geqslant0; \quad ax^2+bx+c > 0\]

Ricordiamo che risolvere una disequazione in un'incognita significa 
determinare l'insieme dei valori da attribuire all'incognita affinché la 
disequazione sia verificata.
Come per le disequazioni di primo grado anche per quelle di secondo grado
dobbiamo studiare il segno del trinomio e poi individuare l'insieme soluzione.

\subsection{Studio del segno di un trinomio di secondo grado}
\label{subsec:diseq_trinomio}
\ind{studio del segno}

Per studiare il segno di un polinomio di secondo grado seguiamo un 
procedimento in due passi:

\begin{enumerate}
\item \textbf{Equazione associata}:
risolvendo l'equazione associata (E.A.) al polinomio troviamo gli zeri della
funzione, cioè i valori della variabile che rendono nullo il polinomio e che
indicano i punti in cui il grafico della funzione incontra l'asse \(x\). 
\ind{zeri di una funzione} \ind{equazione associata E.A.}

A seconda del valore del discriminante l'equazione può avere 
nessuna soluzione reale (\(\Delta<0\)), 
soluzioni reali coincidenti (\(\Delta=0\)), 
soluzioni reali distinte (\(\Delta>0\)).

\item \textbf{Grafico della funzione associata}:
disegnando il grafico della funzione associata al polinomio si individuano 
gli intervalli in cui il polinomio è positivo, cioè il grafico è sopra 
l'asse \(x\) o negativo, cioè il grafico è sotto l'asse \(x\).

Se il coefficiente del termine di secondo grado è minore di zero (\(a < 0\)) 
la parabola ha la concavità rivolta verso il basso, se è maggiore di zero 
(\(a > 0\)) la parabola ha la concavità rivolta verso l'alto. 
\end{enumerate}

Vediamo alcuni esempi di studio del segno di un polinomio di secondo 
grado osservando i diversi casi possibili che dipendono dalla concavità 
della parabola e dall'esistenza di zeri reali della funzione.

\newcommand{\esebasesegno}[6]{% Gli esempi per lo studio del segno.
\begin{esempio}{}{}
Studia il segno del polinomio: \quad \(#1\)

\begin{enumerate}[left=0mm]
\item Equazione Associata al polinomio: \quad \(#1=0\) \\
#2\\
Il polinomio #3.
\item Funzione Associata al polinomio: \quad \(y = #1\)\\
La funzione ha la concavità verso #4.\\
\raisebox{-6mm}{#5}
\end{enumerate}

\vspace{-.8em}
#6
\end{esempio}
}

\esebasesegno{x^2-2x-3}
{\(\tonda{x-3}\tonda{x+1} = 0 \sRarrow x_1=-1;~x_2=+3\)}
{ha due zeri reali distinti}
{l'alto}
{\segniparauz{-1}{+3}}
{La funzione è positiva fino a \(-1\), qui vale \(0\), poi è negativa 
fino \(+3\), qui vale \(0\), poi è positiva.}

\esebasesegno{-x^2+2x+1}
{\(x_{1,~2} = 1 \mp \sqrt{1+1} \sRarrow x_{1,~2} =1 \mp \sqrt{2}\)}
{ha due zeri reali distinti}
{il basso}
{\segniparadz{1-\sqrt{2}}{1+\sqrt{2}}}
{La funzione è negativa fino a \(1-\sqrt{2}\), qui vale \(0\), poi è 
positiva fino \(1+\sqrt{2}\), qui vale \(0\), poi è negativa.}

\esebasesegno{4x^2+4x+1}
{\(\tonda{2x+1}\tonda{2x+1} = 0 \sRarrow x_{1,~2} = -\dfrac{1}{2}\)}
{ha due zeri reali coincidenti}
{l'alto}
{\segniparauzc{-\dfrac{1}{2}}{+}}
{La funzione è positiva su tutto \(\R\) tranne che in \(-\dfrac{1}{2}\), 
dove vale \(0\).}

\esebasesegno{x^2-6x+9}
{\(\tonda{x-3}\tonda{x-3} = 0 \sRarrow x_{1,~2} = +3\)}
{ha due zeri reali coincidenti}
{il basso}
{\segniparadzc{+3}{-}}
{La funzione è negativa su tutto \(\R\) tranne che in \(+3\), 
dove vale \(0\).}

I prossimi due esempi possono trarre in inganno perché l'equazione 
associata al polinomio non ha soluzioni reali, a volte si dice che è 
impossibile. 
Di fronte a cose ``impossibili'' ci viene istintivo lasciar perdere e
abbandonare l'esercizio.
Ma se guardiamo con più attenzione quello che stiamo facendo, il fatto che 
l'equazione non abbia soluzioni reali significa che il polinomio non ha 
zeri, cioè che il grafico della parabola non interseca l'asse \(x\).
Questa non è una situazione impossibile significa che la parabola sta tutta 
sopra (o tutta sotto) all'asse \(x\), quindi che il polinomio è sempre 
positivo (o sempre negativo) per qualunque valore della sua variabile. 

\esebasesegno{x^2+36}
{\(x^2 = -36 \sRarrow \text{Non Ha Soluzioni Reali}\)}
{non ha zeri}
{l'alto}
{\segniparaui{+}}
{La funzione è positiva su tutto \(\R\).}

\esebasesegno{-x^2-x-1}
{\(x_{1,~2} = \dfrac{-1 \mp \sqrt{1-4}}{2} \sRarrow N.H.S.R.\)}
{non ha zeri}
{il basso}
{\segniparadi{-}}
{La funzione è negativa su tutto \(\R\).}

Una volta studiato il segno del polinomio, risolvere una disequazione nella 
forma \(f(x) \lessgtr 0\) consiste nell'aggiungere un ulteriore passo 
piuttosto banale: individuare l'insieme soluzione.

\newcommand{\esebasedis}[4]{% Gli esempi per le disequazioni.
\begin{esempio}{}{}
Risolvi la disequazione: \quad #1

\begin{enumerate}[left=0mm]
\item Equazione Associata: \quad #2 \\
\item #3
\item #4
\end{enumerate}
\end{esempio}
}

\esebasedis{\(x^2+x-2 > 0\)}
{\(x^2+x-2 = 0\) \\ 
\(\tonda{x-1}\tonda{x+2} = 0 \sRarrow x_1 = -2;~x_2 = +1\)}
% \eqdue{\tonda{x-1}\tonda{x+2} = 0} {-2}{+1}}
{\funzioneassociata{\(y = x^2+x-2 \quad \rightarrow\)}
                  {\segniparauz{-2}{+1}}}
{\inssol{\inteaa{4}{-1.5}{+1.5}{-2}{+1}}
       {\(x<-2 \sor x>+1\)}
       {\(\intervaa{-\infty}{-2}~\cup~\intervaa{1}{-\infty}\)}}

\begin{osservazione}{}{}
L'Insieme Soluzione di una disequazione ha diverse rappresentazioni. 
Tutti e tre questi metodi sono importanti e vengono utilizzati in diverse 
situazioni.
\begin{description}[nosep]
\item [Forma Grafica:] \quad 
\raisebox{-3mm}{\inteaa{4}{-1.5}{+1.5}{-2}{+1}}\\
Usata in particolare quando si deve operare sugli insiemi con 
operazioni insiemistiche come l'unione o l'intersezione.
\item [Con operatori relazionali:] \quad
\(x<-2 \sor x>+1\)\\
Usata in particolare nella definizione di insiemi e nei risultati dei 
libri di testo.
\item [Con le parentesi:] \quad 
\(\intervaa{-\infty}{-2}~\cup~\intervaa{1}{-\infty}\)\\
Utile quando vogliamo mettere in evidenza gli estremi degli intervalli.
\end{description}
\medskip
Perciò è importante conoscerli e saperli usare tutti e tre.
\end{osservazione}

\esebasedis{\(x^2-4x+4 \leqslant 0\)}
{\(x^2-4x+4 = 0\) \\ 
\(\tonda{x-2}^2 = 0 \sRarrow x_{1,~2} = +2\)}
{\funzioneassociata{\(y = x^2-4x+4 \quad \rightarrow\)}
                   {\segniparauzc{+2}}}
{\inssol{\intpunto{4}{0}{+2}{blue}}
        {\(x=2\)}
        {\(\lbrace +2 \rbrace\)}}

\esebasedis{\(-x^2+2x-7 > 0\)}
{\(-x^2+2x-7 = 0\) \\ 
\(\dfrac{\Delta}{4}= 1 -7 = -6 \sRarrow \text{N.H.S.R.}\)}
{\funzioneassociata{\(y = -x^2+2x-7 \quad \rightarrow\)}
  {\segniparadi}}
{\inssol{\intv{4}}
  {\(\nexists x \in \R\)}
  {\(\emptyset\)}}

\esebasedis{\(+x^2-3x+10 > 0\)}
{\(+x^2-3x+10 = 0\) \\ 
\(\Delta = 9 -40 = -31 \sRarrow \text{N.H.S.R.}\)}
{\funzioneassociata{\(y = +x^2-3x+10 \quad \rightarrow\)}
  {\segniparaui}}
{\inssol{\intr{4}}
  {\(\intervaa{-\infty}{+\infty}\)}
  {\(\R\)}}

\esebasedis{\(2x^2-5 \leqslant 0\)}
{\(2x^2-5 = 0\) \\ 
\(x^2 =\dfrac{5}{2} \sRarrow x_{1, 2}= \mp \sqrt{\dfrac{5}{2}}\)}
{\funzioneassociata{\(y = 2x^2-5 \quad \rightarrow\)}
  {\segniparauz{-\sqrt{\dfrac{5}{2}}}{+\sqrt{\dfrac{5}{2}}}}}
{\inssol{\inticc{4}{-1.5}{+1.5}{-\sqrt{\frac{5}{2}}}{+\sqrt{\frac{5}{2}}}}
  {\(-\sqrt{\dfrac{5}{2}} \leqslant x \leqslant +\sqrt{\dfrac{5}{2}}\)}
  {\(\intervcc{-\sqrt{\dfrac{5}{2}}}{+\sqrt{\dfrac{5}{2}}}\)}}

\esebasedis{\(-3x^2 + 2x > 0\)}
{\(-3x^2 + 2x=0\) \\ 
\(x \tonda{-3x +2} = 0 \sRarrow x_{1}= 0 \sand x_{2}= \dfrac{2}{3}\)}
{\funzioneassociata{\(y = -3x^2 + 2x \quad \rightarrow\)}
  {\segniparadz{0}{+\dfrac{2}{3}}}}
{\inssol{\intiaa{4}{-1.5}{+1.5}{0}{+\frac{2}{3}}}
  {\(0 < x < +\dfrac{2}{3}\)}
  {\(\intervaa{0}{+\dfrac{2}{3}}\)}}

% \esebasedis{}
% { \\ 
% }
% {\funzioneassociata{}
%   {}}
% {\inssol{}
%   {}
%   {}}

% %%%%%%%%%%%%%%%%%%%%%%%%%%%%%%
% \section{Segno del trinomio a coefficienti letterali}
% Consideriamo il trinomio \(t=kx^2+3x-7\) di secondo grado avente il primo 
% coefficiente dipendente dal parametro \(k\). Come possiamo stabilire il 
% segno 
% di 
% questo trinomio, al variare di \(k\(?
% Sappiamo che stabilire il segno di un trinomio significa determinare i 
% valori 
% reali che attribuiti alla variabile indipendente \(x\) rendono il trinomio 
% positivo, nullo o negativo. Evidentemente per valori reali diversi di \(k\) 
% avremo 
% una diversa disequazione da risolvere; dobbiamo dunque cercare di 
% analizzare 
% come varia il trinomio a seconda dei valori di \(k\) e in seguito studiare 
% il 
% segno del trinomio ottenuto. Questa analisi di situazioni diverse è la 
% \textit{discussione del trinomio a coefficienti parametrici}.
% 
% %%%%%%%%%%%%%%%%%%%%%%%%%%%%%%
% \conclusione Una disequazione di secondo grado si presenta sempre in una 
% delle 
% seguenti forme: \({ax}^2+{bx}+c>0\), \({ax}^2+{bx}+c\geqslant 0\), 
% \({ax}^2+{bx}+c<0\), 
% \({ax}^2+{bx}+c\leqslant 0\) possiamo sempre supporre positivo il primo 
% coefficiente 
% e, 
% anche se incompleta, per l'equazione associata possiamo sempre pensare ai 
% tre 
% casi generati dal segno del discriminante \(\Delta =b^2-4{ac}\). Pertanto 
% avremo:
% 
% %%%%%%%%%%%%%%%%%%%%%%%%%%%%%
% \begin{esempio}{}{}
% Stabilire il segno di \(t=kx^2+3x-7\) al variare di \) k \).
% 
% Prendiamo in considerazione il segno del primo coefficiente e il segno 
% del 
% discriminante dell'equazione associata \(kx^2+3x-7=0\). Il primo 
% coefficiente è 
% maggiore di zero per \(k>0\). Il discriminante \(\Delta =9+28k\) è 
% maggiore 
% di 
% zero 
% per \(k>-\frac 9{28}\). Rappresentiamo la loro reciproca situazione:
% \begin{center}
% %%%%%%%%%%%%%%%%%%%%%%%%%%%%%%
% \begin{threeparttable}
% \begin{tabular}{lllll}
% \toprule
% Delta & \(ax^2+bx+c>0\)& \(ax^2+bx+c\geqslant0\(& \(ax^2+bx+c<0\) & 
% \(ax^2+bx+c\leqslant0\)\\
% \midrule
%  \(\Delta >0^{*}\(& \) x<x_1\vee x>x_2 \) & \) x\leqslant x_1\vee 
% x\geqslant % x_2 % \(& \) 
% x_1<x<x_2 
% \(&\) x_1\leqslant x\leqslant x_2 \(\\
% \(\Delta =0^{**}\)& \(\forall x\in \insR-\{x_1\} \) & \) \forall x\in 
% \insR 
% \(& \) 
% \IS=\emptyset \(&\) x=x_1=x_2 \(\\
% \(\Delta <0^{***}\(&\) \forall x\in \insR \) & \) \forall x\in \insR \(& 
% \) 
% \IS=\emptyset \(&\) \IS=\emptyset \(\\
% \bottomrule
% \end{tabular}
% \begin{tablenotes}
% \item [*] L'equazione associata ha 2 soluzioni reali distinte: 
% \(x=x_1\vee 
% x=x_2\).
% \item [**] L'equazione associata ha 2 soluzioni reali coincidenti: 
% \(x=x_1=x_2\).
% \item [***] L'equazione associata non ha soluzioni reali.
% \end{tablenotes}
% \end{threeparttable}
% \end{center}
% %%%%%%%%%%%%%%%%%%%%%%%%%%%%%%
%  \input{\folder lbr/fig019_seg.pgf}
% % \end{center}
% \vspazio\ovalbox{\risolvii \ref{ese:4.1}, \ref{ese:4.2}, \ref{ese:4.3}, 
% \ref{ese:4.4}, \ref{ese:4.5}, \ref{ese:4.6}}
% 
% \end{esempio}
% 
% %%%%%%%%%%%%%%%%%%%%%%%%%%%%%%

\section{Disequazioni di polinomi scomponibili}
\label{sec:diseq_scomposizione}

A volte potremmo trovarci di fronte al problema di risolvere una
disequazione di grado superiore al secondo.
Potrebbe essere un problema difficile da risolvere, ma, se riusciamo a
scomporre in fattori il polinomio, la difficoltà può scomparire.
\ind{disequazioni di grado superiore}

Partiamo con un esempio commentato.

\begin{esempio}{}{}
Risolvi la disequazione: \quad \(p(x)=-2x^3+3x^2+8x-12 \leqslant 0\)

Il polinomio è di terzo grado. 
Possiamo provare a scomporre in fattori il polinomio in modo da 
ottenere il prodotto di più polinomi di grado inferiore. 
In questo caso ci viene in aiuto il metodo del raccoglimento 
parziale:\\[-4mm]
\begin{align*} \small
-2x^3+3x^2+8x-12 &= x^2(-2x+3)-4(-2x+3) = (x^2-4)(-2x+3)
\end{align*}
Possiamo ora studiare il segno dei fattori poi, applicando la 
regola dei segni di un prodotto, troviamo il segno del polinomio 
iniziale e quindi l'insieme soluzione della disequazione.

\begin{itemize}[left=0mm]
\item segno di \(x^2-4\)\\[-2mm]
\zeripol{\eqduesf{x^2-4 = 0 \sRarrow x^2 = +4}{-2}{+2}}\\
\grafun{\funzass{x^2-4}{\segniparauz{-2}{+2}}}
\item  segno di \(-2x+3\)\\%[-1mm]
\zeripol{\(-2x+3 = 0 \sRarrow -2x = -3 \sRarrow x = \dfrac{3}{2}\)}\\
\grafun{\funzass{-2x+3}{\segnidec{\dfrac{3}{2}}}}
\item Segno del prodotto: \quad %\\[-2em]
\raisebox{-10mm}
{\immagine*[.8]{Grafo dei segni del prodotto di due fattori.}
{\grrsegnoprodottoa}}
\end{itemize}

Quindi i valori di~\(x\) che rendono vera la disequazione, cioè i 
valori che rendono~\(f(x)\) non positivo, sono quelli 
che si trovano tra \(-2\) e \(+\dfrac{3}{2}\), estremi compresi oppure 
che si trovano da \(+2\) compreso in poi.\\
\inssol
  {\solprodottoa}
  {\(-2 \leqslant x \leqslant +\dfrac{3}{2} \sor x \geqslant +2\)}
  {\(\intervcc{-2}{+\dfrac{3}{2}}~\cup~\intervca{+2}{+\infty}\)}
\end{esempio}

\begin{procedura}{}{}
Risolvere le disequazioni di grado superiore:
\begin{enumerate} [noitemsep]
\item scomponi il polinomio di grado \(n\) in fattori di primo e secondo 
grado;
\item studia il segno dei singoli fattori;
\item costruisci la tabella dei segni;
\item applica la regola dei segni per calcolare il segno della funzione 
polinomiale;
\item cerca gli intervalli in cui il polinomio dato assume il segno 
richiesto.
\item rappresenta, con i diversi metodi visti, gli intervalli che 
 risolvono la disequazione.
\end{enumerate}
\end{procedura}

Applichiamo la procedura ai seguenti esempi.

% \newpage %----------------------------------------------

\begin{esempio}{}{}
Data la disequazione 
\[-2x(3-2x)-3x^2\tonda{2-\dfrac{3}{2}x} < 5\tonda{2x^2-\dfrac{3}{10}x}\] 
determina il suo \(\IS\)

Osserviamo che la disequazione proposta è polinomiale di terzo grado; 
esegui i calcoli per portarla alla forma \(p(x)\geqslant 0\): 
% \vspace{3em}
\[-6x +4x^2 -6x^2 +\dfrac{9}{2}x^3 -10x^2 +\dfrac{3}{2}x < 0
\ssRarrow
9x^3 -24x^2 -9x < 0\]
Scomponendo in fattori si ottiene: \(3x \cdot (3x^2-8x-3) < 0\). 

A questo punto possiamo studiare il segno dei fattori e, applicando la 
regola dei segni di un prodotto, trovare il segno del polinomio iniziale 
e risolvere la disequazione.

Studiamo il segno di ogni singolo fattore:
\begin{itemize}[left=0mm]
\item segno di \(3x^2-8x-3\)\\[-2mm]
\hspace*{-3mm}\zeripol{
\eqduesf{x_{1,~2} = \dfrac{+4 \mp \sqrt{16 +9}}{3}}
{-\dfrac{1}{3}}{+3}}\\
\hspace*{-3mm}\grafun{\funzass{3x^2-8x-3}{\segniparauz{-\dfrac{1}{3}}{+3}}}
\item  segno di \(3x\)\\%[-1mm]
\zeripol{\(3x = 0 \sRarrow x = 0\)}\\
\grafun{\funzass{3x}{\segnicre{0}}}
\item Segno del prodotto: \quad %\\[-2em]
\raisebox{-10mm}
{\immagine*[.8]{Grafo dei segni del prodotto di due fattori.}
{\grrsegnoprodottob}}
\end{itemize}

Quindi i valori di~\(x\) che rendono vera la disequazione, cioè i 
valori che rendono~\(f(x)\) non positivo, sono quelli 
che si trovano a sinistra di~\(-2\) oppure che si trovano a destra 
di~\(+3\).\\
\inssol
   {\solprodottob}
   {\(x < -\dfrac{1}{3} \sor 0 < x < +3\)}
   {\(\intervaa{+\infty}{-\dfrac{1}{3}}~\cup~\intervaa{0}{+3}\)}
\end{esempio}

\begin{esempio}{}{}
Un numero è tale che aggiungendo al suo cubo il suo quadruplo si ottiene 
un numero maggiore del triplo del suo quadrato aumentato di~\(12\). Determina 
l'insieme soluzione del problema.

La richiesta del problema implica la ricerca dell'Insieme Soluzione della 
disequazione \(x^3+4x>3x^2+12\), di terzo grado nella variabile~\(x\). 
Scriviamo la disequazione in forma canonica, applicando i principi di 
equivalenza: 
\[x^3-3x^2+4x-12>0\] 
Si tratta di una disequazione polinomiale di terzo grado.

Procediamo nella scomposizione in fattori del polinomio:
\[p(x)=x^3-3x^2+4x-12 = x^2 \tonda{x-3}+4 \tonda{x-3} =
      \tonda{x^2+4} \tonda{x-3}\]
Ora possiamo determinare il segno dei vari fattori, ma se siamo 
sufficientemente acuti e pigri possiamo osservare che il primo fattore sarà 
sempre positivo (perché?) quindi il segno del polinomio è equivalente al 
segno del secondo fattore che si determina molto rapidamente\dots
\inssol
  {\disegno{\assex{0}{6}{0} \ray{0}{6}{3}{+3}{white}}}
  {\(x < +3\)}
  {\(\intervaa{+3}{+\infty}\)}
\end{esempio}

\begin{esempio}{}{}
Risolvi la disequazione: \(1-81x^4<0\)

Il binomio a primo membro può essere scomposto in fattori di secondo grado  
in due passi:
\[1-81x^4 = \tonda{1-9x^2}\tonda{1+9x^2}=
  \tonda{-9x^2+1}\tonda{9x^2+1}\]
Anche qui possiamo usare un po' di furbizia e calcolare la soluzione molto 
rapidamente.\\
\inssol
  {\disegno{\assex{0}{6}{0} 
   \inti{0}{2}{4}{-\dfrac{1}{3}}{+\dfrac{1}{3}}{white}{white}}}
  {\(-\dfrac{1}{3} < x < +\dfrac{1}{3}\)}
  {\(\intervaa{-\dfrac{1}{3}}{+\dfrac{1}{3}}\)}
\end{esempio}

\begin{esempio}{}{}
Risolvere la disequazione: \(x^4-4x^2-45 \geqslant 0\)

Il trinomio al primo membro è di quarto grado, ma applicando la 
sostituzione: 
\(x^2=z\)

può essere ricondotto al trinomio di secondo grado:
\(z^2-4z-45 \geqslant 0\)

la cui scomposizione in fattori risulta:
\(\tonda{z-9}\tonda{z+5}\)

quindi la disequazione assegnata diventa: 
\((x^2-9)(x^2+5) \geqslant 0\)

Completa la soluzione della disequazione.\\
\inssol
  {\disegno{\assex{0}{6}{0} 
   \ray{0}{-0.3}{2}{-3}{blue} \ray{0}{6}{4}{+3}{blue}}}
   {\(x \leqslant -3 \sor x \geqslant +3\)}
   {\(\intervac{+\infty}{-3}~\cup~\intervca{+3}{+\infty}\)}
\end{esempio}

% \vspazio\ovalbox{\risolvii \ref{ese:4.32}, \ref{ese:4.33}, \ref{ese:4.34}, 
% \ref{ese:4.35}, \ref{ese:4.36}, \ref{ese:4.37}, \ref{ese:4.38}, 
% \ref{ese:4.39}, 
% \ref{ese:4.40}, \ref{ese:4.41}, \ref{ese:4.42}, \ref{ese:4.43}, \ref 
% {ese:4.44},}
% 
% \vspazio\ovalbox{\ref{ese:4.45}, \ref{ese:4.46}, \ref{ese:4.47}, 
% \ref{ese:4.48}, 
% \ref{ese:4.49}, \ref{ese:4.50}, \ref{ese:4.51}, \ref{ese:4.52}, 
% \ref{ese:4.53}, 
% \ref{ese:4.54}, \ref{ese:4.55}, \ref{ese:4.56}, \ref{ese:4.57}}

% % \newpage %----------------------------------------------

\section{Disequazioni fratte}
\label{sec:diseq_fratte}
\ind{disequazione fratta}

Ricordiamo che una disequazione è \emph{frazionaria} o \emph{fratta} quando 
il suo denominatore contiene l'incognita. 

% \begin{osservazione}{}{}
% \begin{itemize} [nosep]
% \item studiare il segno di una espressione letterale significa stabilire in 
% quale insieme si trovano i valori della variabile che rendono 
% l'espressione negativa, nulla o positiva;
% \item ogni espressione contenente operazioni tra frazioni algebriche ha in 
% generale come risultato una frazione algebrica;
% \item nello studiare il segno terremo conto anche del campo di esistenza 
% della funzione.
% \end{itemize}
% \end{osservazione}

Per risolvere una disequazione fratta si può applicare la seguente procedura 
quasi identica a quella che si utilizza per le funzioni scomponibili in 
fattori.

\begin{procedura}{}{}
Soluzione di una disequazione frazionaria:
\begin{enumerate} [noitemsep]
\item trasporta tutti i termini al primo membro e scrivilo sotto forma di 
 un'unica frazione:\\
\[\frac{N(x)}{D(x)} \lessgtr 0\]
\item scomponi in fattori numeratore e denominatore;
\item studia il segno di ogni fattore;
\item costruisci la tabella dei segni, 
\emph{ricordando di indicare con una crocetta gli zeri del denominatore}
perché gli zeri del denominatore rendono \emph{non definita} la frazione;
\item applica la regola dei segni per calcolare il segno della funzione 
razionale;
\item individua gli intervalli in cui la frazione assume il segno richiesto;
\item rappresenta, con i diversi metodi visti, gli intervalli che 
 risolvono la disequazione.
\end{enumerate}
\end{procedura}

Vediamo attraverso alcuni esempi come procedere.

\begin{esempio}{}{}
 \[\frac{2}{x-7} \geqslant \frac{3}{x+4}\]
Per prima cosa dobbiamo scrivere la disequazione in forma normale:
{\small
 \[\frac{2}{x-7} - \frac{3}{x+4} \geqslant 0 \Rightarrow 
   \frac{2(x+4) - 3(x-7)}{(x-7)(x+4)} = 
   \frac{2 x +8 -3x +21}{(x-7)(x+4)} =
   \frac{-x +29}{x^2-3x-28} \geqslant 0\]
} % end small
A questo punto possiamo studiare il segno dei fattori e, applicando la regola 
dei segni di un prodotto, trovare il segno della disequazione fratta.

Studiamo il segno di ogni singolo fattore e della frazione:
\begin{itemize}[left=0mm]
\item  segno di \(-x +29\)\\%[-1mm]
\zeripol{\(-x +29=0 \sRarrow x=29\)}\\
\grafun{\funzass{-x +29}{\segnidec{+29}}}

\item segno di \(x^2-3x-28\)\\[-2mm]
\hspace*{-3mm}\zeripol{
\eqduesf{x_{1,~2} = \dfrac{+3 \mp \sqrt{9 +112}}{2} = \dfrac{+3 \mp 11}{2}}
  {-4}{+7}}\\
\hspace*{-3mm}\grafun{\funzass{x^2-3x-28}{\segniparauz{-4}{+7}}}

\item Segno del quoziente: \quad %\\[-2em]
\raisebox{-10mm}
{\immagine*[.8]{Grafo dei segni del quoziente di due fattori.}
{\segnofrazionea}}
\end{itemize}

Quindi i valori di~\(x\) che rendono vera la disequazione, cioè 
i valori  che rendono~\(f(x)\) non negativo, sono quelli 
che si trovano a sinistra di~\(-4\) oppure i valori tra \(+7\) e \(+29\) 
\(+7\) escluso e \(+29\) compreso.\\
\inssol
  {\solfrazionea}
  {\(x < -4 \sor 7 < x \leqslant 29\)}
  {\(\intervaa{-\infty}{-4}~\cup~\intervac{7}{29}\)}
\end{esempio}

\begin{esempio}{}{}
Data l'espressione 
\[E = \dfrac{4}{4x^2-1} +\dfrac {1}{2x+1} +\dfrac{x}{1-2x}\]
determinarne il segno, al variare di \(x\) in \(\R\).

Dopo aver osservato che il denominatore comune è:
\(4x^2-1 =\tonda{2x-1}\tonda{2x+1}\), scriviamo 
l'espressione sotto forma di un'unica frazione:
\[E=\dfrac{4+2x-1 - x \tonda{2x+1}}{\tonda{2x-1}\tonda{2x+1}} =
\dfrac{-2x^2+x+3}{4x^2-1}\]

A questo punto possiamo studiarne il segno.
\begin{itemize}[left=0mm]
\item segno di \(-2x^2 +x +3\)\\[-2mm]
\hspace*{-3mm}\zeripol{
\eqduesf{x_{1,~2} = \dfrac{-1 \mp \sqrt{1 +24}}{-4} = \dfrac{-1 \mp 5}{-4}}
  {+\dfrac{3}{2}}{-1}}\\
\hspace*{-3mm}\grafun{\funzass{-2x^2+x+3}{\segniparadz{-1}{+\dfrac{3}{2}}}}

\item segno di \(4x^2 -1\)\\[-2mm]
\hspace*{-3mm}\zeripol{
\eqduesf{x^2 = \dfrac{1}{4} \sRarrow}
  {-\dfrac{1}{2}}{+\dfrac{1}{2}}}\\
\hspace*{-3mm}
\grafun{\funzass{4x^2-1}{\segniparauz{-\dfrac{1}{2}}{+\dfrac{1}{2}}}}

\item Segno del quoziente: \quad %\\[-2em]
\raisebox{-10mm}
{\immagine*[.8]{Grafo dei segni del quoziente di due fattori.}
{\segnofrazioneb}}
\end{itemize}

% Quindi i valori di~\(x\) che rendono vera la disequazione, cioè 
% i valori  che rendono~\(f(x)\) non negativo, sono quelli 
% che si trovano a sinistra di~\(-4\) oppure i valori tra \(+7\) e \(+29\) 
% \(+7\) escluso e \(+29\) compreso.\\
% \inssol
%   {\solfrazionea}
%   {\(x < -4 \sor 7 < x \leqslant 29\)}
%   {\(\intervaa{-\infty}{-4}~\cup~\intervac{7}{29}\)}
\end{esempio}

\begin{esempio}{}{}
Determina l'Insieme Soluzione della disequazione fratta:\\ 
\(3-\dfrac{1}{2x+1} \geqslant \dfrac{1}{1-x}\)

\medskip
Dopo aver osservato che il denominatore comune è:
\[\tonda{2x+1}\tonda{-x+1} = -2x^2 +2x -x +1 = -2x^2 +x +1\]
scrivi l'espressione sotto forma di un'unica frazione:
\[3-\dfrac{1}{2x+1} - \dfrac{1}{1-x} =
  \dfrac{3 \tonda{-2x^2 +x +1} - \tonda{-x+1} - \tonda{2x+1}}
        {\tonda{2x+1}\tonda{-x+1}} =\]
\[\dfrac{-6x^2 +3x +3 +x-1 -2x-1}
        {\tonda{2x+1}\tonda{-x+1}} = 
  \dfrac{-6x^2 +2x +1}
        {-2x^2 +x +1} \geqslant 0\]
A questo punto possiamo studiarne il segno.

\begin{itemize}[left=0mm]
\item segno di \(-6x^2 +2x +1\)\\[-2mm]
\hspace*{-3mm}\zeripol{
\eqduesf{x_{1,~2} = \dfrac{-1 \mp \sqrt{1 +6}}{-4} = 
\dfrac{-1 \mp \sqrt{7}}{-6}}
  {\dfrac{1 - \sqrt{7}}{6}}{\dfrac{1 + \sqrt{7}}{6}}}\\
\hspace*{-3mm}
\grafun{\funzass{-6x^2 +2x +1}
{\segniparadz{\dfrac{1 - \sqrt{7}}{6}}{\dfrac{1 + \sqrt{7}}{6}}}}

\item segno di \(-2x^2 +x +1\)\\[-2mm]
\hspace*{-3mm}\zeripol{
\eqduesf{x_{1,~2} = \dfrac{-1 \mp \sqrt{1 +8}}{-4} = 
\dfrac{-1 \mp 3}{-4}}
  {+1}{-\dfrac{1}{2}}}\\
\hspace*{-3mm}
\grafun{\funzass{4x^2-1}{\segniparadz{{-\dfrac{1}{2}}}{+1}}}

\item Segno del quoziente: \quad %\\[-2em]
\raisebox{-10mm}
{\immagine*[.8]{Grafo dei segni del quoziente di due fattori.}
{\segnofrazionec}}
\end{itemize}

Quindi i valori di~\(x\) che rendono vera la disequazione, cioè 
i valori  che rendono~\(f(x)\) non negativo, sono quelli 
che si trovano a sinistra di~\(-\dfrac{1}{2}\) oppure 
i valori tra \(\dfrac{1 - \sqrt{7}}{6}\) e \(\dfrac{1 + \sqrt{7}}{6}\) 
estremi compresi o
i valori oltre \(+1\).\\
\inssol
  {\solfrazionec}
  {\(x < -\dfrac{1}{2} \sor 
  \dfrac{1 - \sqrt{7}}{6} \leqslant x \leqslant \dfrac{1 + \sqrt{7}}{6} 
  \sor x > +1\)}
  {\(\intervaa{-\infty}{-\dfrac{1}{2}}~\cup~
  \intervcc{\dfrac{1 - \sqrt{7}}{6}}{\dfrac{1 + \sqrt{7}}{6}}~\cup~
  \intervaa{+1}{+\infty}\)}
%  Studia il segno di ogni singolo fattore:
% 
% \begin{itemize}
% 
%  \item  segno di \(-6x^2 +2x +1\)\\
%  \segnofatt
%    {\(-6x^2 +2x +1=0 \sRarrow\)}
%    {\(y=-6x^2 +2x +1\)}
%    {[parabola concavità verso il basso, }
%    {\segniparadz{\frac{1-\sqrt 7}{6}}{\frac{1+\sqrt 7}{6}}}
%  \item segno di \(-2x^2 +x +1\)\\
%  \segnofatt
%    {\(-2x^2 +x +1=0 \sRarrow \) \\
%   \(x_{1,2}=\mp\dfrac{1}{2}\)}
%    {\(y=-2x^2 +x +1\)}
%    {parabola concavità verso il basso}
%    {\segniparadz{}{}}
% \end{itemize}
% 
% Riporta i segni del numeratore e del denominatore e calcola
% il segno della frazione:
% 
% \begin{inaccessibleblock}[Grafo con i segni]
% \vspace{2em}
%   \begin{center}
%   \segnofrazionec
%   \end{center}
% \end{inaccessibleblock}
% 
%  \insiemesoluzione{\intv{8}}{}{}
\end{esempio}

% \vspazio\ovalbox{\risolvii \ref{ese:4.58}, \ref{ese:4.59}, \ref{ese:4.60}, 
% \ref{ese:4.61}, \ref{ese:4.62}, \ref{ese:4.63}, \ref{ese:4.64}, 
% \ref{ese:4.65}, 
% \ref{ese:4.66}, \ref{ese:4.67}, \ref{ese:4.68}, \ref{ese:4.69}, \ref 
% {ese:4.70},}
% 
% \vspazio\ovalbox{\ref{ese:4.71}, \ref{ese:4.72}, \ref{ese:4.47}, 
% \ref{ese:4.73}, 
% \ref{ese:4.74}}

\section{Sistemi di disequazioni}
\label{sec:diseq_sistemi}

Ricordiamo che risolvere un sistema di disequazioni significa trovare 
l'insieme dei numeri reali che sono le soluzioni comuni alle disequazioni 
che lo compongono, cioè l'intersezione delle soluzioni delle singole 
disequazioni. 
\ind{sistema di disequazioni}

\begin{definizione}{}{}
Indicate con \(D_{1}, D_{2}, \ldots, D_n\) le disequazioni che formano 
il sistema e \(\IS_{1}, \IS_{2}, \ldots,\IS_n\) i rispettivi insieme 
soluzione, la soluzione del sistema indicata con \(\IS\) è data da 
\[\IS=\IS_1 ~\cap~ \IS_2 ~\cap~ \ldots ~\cap~ \IS_n\]
\end{definizione}

Per risolvere un sistema di disequazioni, bisogna quindi risolvere ogni 
singola disequazione e poi operare l'intersezione fra questi insiemi, cioè 
trovare i valori che hanno in comune.

\begin{osservazione}{}{}
Nelle singole disequazioni non occorre scrivere le soluzioni nei diversi 
metodi, basta scriverla in forma grafica.
\end{osservazione}

% % \newpage %---------------------------------

\begin{esempio}{}{}
 
Risolvi il seguente sistema di disequazioni:
\[\sistema{
    2x^3+3x^2-18x-27 > 0 & \slarrow D_1\\
    \dfrac{-x^2 -2x+15}{4x^2-9} \leqslant 0 & \slarrow D_2
  }
\]
Se scomponiamo in fattori i diversi polinomi ci risulterà più facile 
risolvere le equazioni associate.
\[\sistema{
\tonda{x^2-9} \tonda{2x+3} > 0 & \slarrow D_1\\[3mm]
\dfrac{-\tonda{x-3}\tonda{x+5}}{\tonda{2x-3}\tonda{2x+3}} \leqslant 0 & 
\slarrow D_2
}\]
Risolviamo una disequazione alla volta studiando i segni dei polinomi:

\begin{itemize}[left=0mm]
\item 
\(D_1: \quad 2x^3+3x^2-18x-27 > 0 \ssLRarrow \tonda{x^2-9}\tonda{2x+3} > 0\)

\affiancati{.44}{.54}{
Tenendo conto degli zeri dei polinomi 
\(\graffa{-\frac{3}{2}}\) e \(\graffa{-3;~+3}\), della pendenza 
della retta associata e della concavità della parabola associata possiamo 
risolvere la disequazione usando il grafo a fianco.
}{
\immagine[.8]{Grafo dei segni per la prima disequazione.}{\segnosistemaaa}
}

\item 
\(D_1: \quad \dfrac{-x^2 -2x+15}{4x^2-9} \leqslant 0
\ssLRarrow
\dfrac{-\tonda{x-3}\tonda{x+5}}{\tonda{2x-3}\tonda{2x+3}} \leqslant 0 
\)

Dato che abbiamo già scomposto in fattori, possiamo 
riconoscere che gli zeri del numeratore sono:
\(\graffa{-5;~3}\) 
e gli zeri del denominatore sono:
\(\graffa{-\frac{3}{2};~+\frac{3}{2}}\) 

\affiancati{.29}{.69}{
La concavità della prima parabola è verso il basso, 
la concavità della seconda è verso l'alto otteniamo così il grafo:
}{
\immagine[.8]{Grafo dei segni per la seconda disequazione.}{\segnosistemaab}
}
\end{itemize}

% \begin{minipage}{.49\textwidth}
Ora riportiamo su un grafo le soluzioni delle due disequazioni. Il grafo deve
anche avere un asse~x sul quale indicare l'intersezione dei due insiemi 
soluzione:
\immagine{Grafo degli insiemi soluzione delle due disequazioni e della loro 
intersezione.} 
{\grrsistemaa}

Possiamo ora rappresentare l'insieme soluzione con gli altri metodi che 
conosciamo:

\begin{minipage}{.49 \linewidth}
\[x \leqslant -5 \sor -\dfrac{3}{2} < x < +\dfrac{3}{2}\]
\end{minipage}
\hfill
\begin{minipage}{.49 \linewidth}
\[\intervac{-\infty}{-5}~\cup~\intervaa{-\dfrac{3}{2}}{+\dfrac{3}{2}}\]
\end{minipage}
\end{esempio}

\begin{esempio}{}{}

In quest'altro esempio vogliamo calcolare l'insieme di esistenza della 
seguente funzione:

\[y = \sqrt{\dfrac{-x^2 +3x +21}{x-7}} + \sqrt{\dfrac{-x -5}{-x^2 +16}}\]

\medskip
Per poter calcolare una radice quadrata nei numeri reali, il radicando deve 
essere non negativo. 
Ci sono due radici quindi abbiamo due condizioni:

\medskip
\affiancati{.49}{.49}{
Le due condizioni:
\[\sistema{
    \dfrac{-x^2 +3x +21}{x-7} \geqslant 0 & C_1\\
    \dfrac{-x -5}{-x^2 +16} \geqslant 0 & C_2
  }
\]
}{
Lo stesso sistema con i polinomi scomposti in fattori:
\[\sistema{
    \dfrac{-\tonda{x+4} \tonda{x-7}}{x-7} \geqslant 0 \\
    \dfrac{-x -5}{-\tonda{x -4} \tonda{x+4}} \geqslant 0
  }
\]
}

Risolviamo una disequazione alla volta studiando i segni dei polinomi:

\begin{itemize}[left=0mm]
\item \(C_1: \quad \dfrac{-x^2 +3x +21}{x-7} \geqslant 0\)

Osservando i vari polinomi scomposti in fattori possiamo ricavare gli zeri 
del numeratore: \(\graffa{-4:~+7}\) e del denominatore: \(\graffa{+7}\)

\affiancati{.49}{.49}{
Osservando il trinomio e il binomio possiamo vedere che la concavità della 
parabola al numeratore è verso il basso e la retta al 
denominatore è crescente. Possiamo così ricavare il grafo a fianco.
}{
\immagine[.8]{Grafo dei segni e soluzione del primo sistema.}
{\segnosistemaba} 
}

\item \(C_2: \quad \dfrac{-x -5}{-x^2 +16} \geqslant 0\)

\affiancati{.49}{.49}{
Anche in questo caso, dato che abbiamo già scomposto in fattori possiamo 
riconoscere che gli zeri del numeratore sono:
\(\graffa{-5}\) 
e gli zeri del denominatore sono:
\(\graffa{-4;~+4}\) 

La retta, al numeratore, è decrescente e 
la concavità della parabola, al denominatore, è verso il basso, quindi 
otteniamo il grafo:
}{
\immagine[.8]{Grafo dei segni e soluzione del secondo sistema.}
{\segnosistemabb}
}

\end{itemize}

\begin{minipage}{.49\textwidth}
Ora riportiamo su un grafo le soluzioni delle due disequazioni.Il grafo deve 
anche avere un asse~x sul quale indicare l'intersezione dei due insiemi 
soluzione:
\end{minipage}
\hfill
\begin{minipage}{.49\textwidth}
\begin{center} \grrsistemab \end{center}
\end{minipage}
% \vspace{-.5em}
% \vspace{-.5em}
L'espressione precedente dà risultati reali solo per valori di~x compresi 
tra~\(-5\) e \(-4\) con il primo estremo incluso e il secondo escluso:\\ 
\(-5 \leqslant x < -4 \qquad \intervca{-5}{-4}\)

\end{esempio}
